\scanpage{14}

Quant : $R,R'$, ce sont respectivement les images de $E=\operatorname{Rep}$ dans $S\times A$ et $A\times F$ ....

\uline{Corollaire} \deleted{Les autom. de} gpe $\operatorname{Aut}(\Pi)$ s'identifie à $\operatorname{Aut}_{\unsure{G|\mathcal G|\mathfrak G|\mathfrak S}_2\text{-ens}}(\operatorname{Rep}(\Pi))$.

\uline{Corollaire} Conditions équivalentes

a) Deux sommets de $\Pi$ sont \inserted{tjrs} reliés par une chaine d'\unsure{arê\textsuperscript{tes}|aires|arêtes}

a$'$) Deux faces de $\Pi$ sont \inserted{tjrs} reliées par une chaine d'arêtes

\[
 \text{b)}\quad\Pi\simeq\Pi'\amalg\Pi''\ \Longrightarrow\ \Pi'=\varnothing\ \text{ou}\ \Pi''=\varnothing.
\]

c) \deleted{\unsure{le}} $G_2$ opère \uline{transitivement} dans $E$\inserted{$=\operatorname{Rep}(\Pi)$}. (i.e. -- si $\Pi\ne\varnothing$ i.e. $E\ne\varnothing$ \deleted{\unsure{\}}} -- $E$ est un \uline{espace homogène} sous $G_2$)

\deleted{Corollaire} \uline{Déf} On dit alors que $\Pi$ est connexe.

\uline{Corollaire} Si $\Pi$ est connexe, $\operatorname{Aut}\Pi$ opère librement dans $\operatorname{Rep}(\Pi)$, i.e. si $g\in\operatorname{Aut}(\Pi)$, $r\in\operatorname{Rep}(\Pi)$, alors

\[
 g.r=r\quad\Longrightarrow\quad g=\mathrm{id}
\]

En effet, un autom d'un espace homogène qui fixe un pt, \deleted{fixe l'espace homogène.} \inserted{est l'identité}.

% archive sheet number: 13

